\documentclass[12pt]{exam}

\include{mypreamble}

\begin{document}

\title{Problem Set 8}
\date{Due: April 13th, 5pm}
\maketitle



\section{Formal proofs}
For each problem in this section, you must give a complete and rigorous proof. Your arguments should be clear, logically ordered, and written in full sentences.

In particular, state all assumptions explicitly and define all notation. Justify every nontrivial step. 

\begin{questions}
\question[1] Let $p$ be an odd prime. Compute the value of the Legendre symbol $\left(\frac{-1}{p}\right)$ as a formula depending on $p \pmod{4}$. Prove your formula is correct. 

\question[1] Let $a, b \in \mathbb{R}^+$. Prove that the number of integer lattice points $(x, y)$ contained strictly inside the the rectangle with corners $(0,0)$ and $(a,b)$ is given by $\lfloor a \rfloor \lfloor b \rfloor$ for the number of points in the region $1 \le x \le a, 1 \le y \le b$ when $(a, b)$ is not an integer lattice point.

\question[1] Let $p$ be an odd prime and let $a, b$ be integers. Prove using the definition that the Legendre symbol modulo $p$ is multiplicative, that is,
$$\left(\frac{ab}{p}\right) = \left(\frac{a}{p}\right)\left(\frac{b}{p}\right).$$
(To receive any points, you must provide a complete proof directly from the definition, and not using any other properties of the Legendre symbol such as Euler's Criterion or Gauss's Lemma.)

\section{Demonstrations}
\fullwidth{For problems in this section, you still need to give complete mathematical reasoning to support your answers. I should be able to follow and understand your work, but it doesn't have to be organized into a formal proof. }

\question[1] Let $p$ be an odd prime. Using Gauss's Lemma, compute the value of the Legendre symbol $\left(\frac{2}{p}\right)$. (Your answer should be in terms of $p \pmod{8}$.)

\question[1] The \emph{Jacobi symbol} generalizes the Legendre symbol to odd composite moduli. Let $n > 1$ be an odd integer with prime factorization $n = p_1^{a_1} p_2^{a_2} \cdots p_k^{a_k}$, and let $a$ be any integer. The Jacobi symbol $\left(\frac{a}{n}\right)$ is defined as
$$\left(\frac{a}{n}\right) = \left(\frac{a}{p_1}\right)^{a_1} \left(\frac{a}{p_2}\right)^{a_2} \cdots \left(\frac{a}{p_k}\right)^{a_k},$$
where the factors on the right are Legendre symbols. By convention, $\left(\frac{a}{1}\right) = 1$.

\begin{parts}
\part Show that the Jacobi symbol is multiplicative in the top, that is, $\left(\frac{ab}{n}\right) = \left(\frac{a}{n}\right)\left(\frac{b}{n}\right)$.

\part Show that the Jacobi symbol is multiplicative in the bottom, that is, if $m$ and $n$ are odd positive integers, then $\left(\frac{a}{mn}\right) = \left(\frac{a}{m}\right)\left(\frac{a}{n}\right)$.

\part Show that $\left(\frac{a}{n}\right) = 0$ if and only if $\gcd(a, n) > 1$.
\end{parts} 


\question[1] Use the Law of Quadratic Reciprocity and properties of the Jacobi symbol to compute the following. Show all steps in your computation.
\begin{parts}
    \part $\left(\frac{105}{48379}\right)$
    \part $\left(\frac{219}{383}\right)$
    \part $\left(\frac{6244334}{567617}\right)$
\end{parts}

\question[1] For $p = 11$, compute the Legendre symbol $\left(\frac{a}{11}\right)$ for all $a \in \{2, 5, 8, 9\}$ using three different methods:
\begin{parts}
    \part By definition (checking if $a$ is a quadratic residue by finding all squares modulo 11).
    \part Using Euler's Criterion.
    \part Using Gauss's Lemma (compute the set $S(a, 11)$ for each $a$ and compute $\mu(a,11)$).
\end{parts}



\end{questions}
\end{document}
